\documentclass[UTF8,zihao=-4]{ctexart}

\usepackage[b5paper,margin=17mm,top=18mm,bottom=18mm]{geometry}
\usepackage{xcolor}
\usepackage{amsmath}
\usepackage{enumitem}
\usepackage{ulem}
\usepackage{setspace}

\definecolor{workbookgray}{gray}{0.35}
\setlength{\parindent}{2em}
\setlength{\parskip}{0.35em}
\setstretch{1.35}
\pagestyle{empty}

\newcommand{\question}[2]{%
  \par\noindent\textbf{#1}\quad #2\par
}

\newcommand{\marked}[1]{%
  \ensuremath{\overset{\scriptstyle\cdot}{\text{#1}}}%
}

\begin{document}

\begin{center}
  {\zihao{3}\bfseries 语文、政治习题册模板样章}
  \\
  {\small\color{workbookgray}公开模板示例，不含真实题库}
\end{center}

\section*{第一章\hspace{1em}题目组件示例}

\noindent\textbf{使用说明：}本页只展示题目、括号答题位置和基本层级。正式项目中，内容必须经过来源登记与教师审核。

\question{1．}{请在括号内填写答案（\hspace{3em}）。}

\question{2．}{下列句子中，\marked{加}\marked{点}词语的意义与现代汉语不同的一项是（\hspace{3em}）。}

\question{3．}{请结合上下文，判断\uline{划线句子}的表达作用（\hspace{7em}）。}

\vspace{1em}
\noindent\rule{\linewidth}{0.4pt}

\begin{center}
  {\small\color{workbookgray}后续版本将在此基础上加入章节、题型、教师版和答案组件。}
\end{center}

\end{document}

